\documentclass[12pt,a4paper]{article}
\usepackage{fontawesome5}
\providecommand{\faPencilAlt}{\faIcon{pencil-alt}}
\usepackage{styles/couleurs}
\usepackage{styles/mise_en_page}
\usepackage{styles/boites}
\usepackage{styles/maths}
\usepackage{styles/figures}
\usepackage{enumitem}
\usepackage{array}
\usepackage{tabularx}
\usepackage{amsmath}
\usepackage{amssymb}
\usepackage{tikz}
\usepackage{xparse}
\usetikzlibrary{calc,patterns,arrows.meta}
\setlist[itemize]{leftmargin=1.4em,itemsep=0.15em,topsep=0.2em}
\setlist[enumerate]{leftmargin=*,itemsep=0.25em,topsep=0.25em}
\newcommand{\lignepointillee}{\noindent\dotfill\par\vspace{0.45em}}
\newcommand{\lignesreponse}[1]{%
\begin{tcolorbox}[enhanced,colback=white,colframe=NeutreC!40,boxrule=0.5pt,
arc=1.5mm,left=3mm,right=3mm,top=2mm,bottom=1mm,before skip=2pt,after skip=2pt]
\foreach \i in {1,...,#1}{\lignepointillee}
\end{tcolorbox}}
\newcommand{\titreseq}[1]{%
\begin{center}
{\sffamily\LARGE\bfseries\color{TitreC}Séquence 4 -- Calcul littéral}\\[0.2em]
{\sffamily\large\bfseries\color{SavoirC}#1}
\end{center}}
\begin{document}

\titreseq{Automatismes -- Séries courtes}

\remarque{Chaque série est prévue pour environ cinq minutes.}

\section*{Série 1 -- Traduire}

\begin{enumerate}
  \item Le triple de \(x\), augmenté de \(5\).
  \item La somme de \(7\) et du produit de \(4\) par \(x\).
  \item Le produit de \(3\) par la somme de \(x\) et de \(2\).
  \item Le carré de \(x\), diminué de \(6\).
\end{enumerate}

\section*{Série 2 -- Calculer une valeur}

Pour \(x=-2\), calcule :

\[
3x+5 \qquad x^2-4 \qquad 2(x+7) \qquad 5-4x.
\]

\section*{Série 3 -- Simplifier l'écriture}

\[
4\times x \qquad x\times7 \qquad x\times x
\qquad 1\times x \qquad 3\times(x+2).
\]

\section*{Série 4 -- Réduire}

\[
3x+5x \qquad x+x+4 \qquad 7x-2x
\qquad 4x+3+2x+5.
\]

\section*{Série 5 -- Réduire avec des relatifs}

\[
-3x+7x \qquad 5x-8x \qquad -4x+6+2x-9
\qquad 3x^2+2x-5x^2.
\]

\section*{Série 6 -- Développer}

\[
3(x+5) \qquad 4(2x-3) \qquad -2(x+7)
\qquad 5(3x-1).
\]

\section*{Série 7 -- Développer puis réduire}

\[
3(x+2)+4x
\qquad
5(2x-1)-3x
\qquad
-2(x+4)+7x.
\]

\section*{Série 8 -- Vrai ou faux ?}

Indique si l'égalité est vraie pour toute valeur de \(x\).
Justifie par réduction ou donne un contre-exemple.

\[
3(x+2)=3x+6
\]

\[
4x+5x=9x
\]

\[
2x+3=5x
\]

\[
5(x-1)+2x=7x-5.
\]

\end{document}
